\section{Summary and Takeaways}

Deep learning has a longer and stranger history than most people realize.

It began with Frank Rosenblatt's Perceptron in 1958 --- a single layer of adjustable weights and a dream of machine intelligence.
It survived two AI winters, decades of skepticism, and the rise and fall of competing techniques.
Backpropagation gave multi-layer networks the ability to learn complex patterns.
LSTMs gave networks a smarter memory for sequences.
GPUs and ImageNet gave the field the scale it had always needed.

Then in 2012, AlexNet announced to the world that deep learning was back --- and this time it meant business.
The decade that followed saw it conquer image recognition, then speech, then translation.

The crucial turn came in 2017, when the Transformer architecture replaced sequential processing with parallel self-attention.
Suddenly, language models could be trained on enormous datasets.
Scaling laws emerged: more data, more parameters, more compute reliably produced better models.

That trajectory led directly to GPT, ChatGPT, Claude, and DeepSeek --- models that can write, reason, code, and converse in ways that felt like science fiction a decade ago.

And through all of it, the foundation has not changed.
Neurons. Weights. Layers. Gradient descent. Backpropagation.

The tools that Frank Rosenblatt first sketched in 1958, refined over sixty years of research, now run on thousands of specialized chips and power some of the most consequential technology ever built.

\textbf{Deep learning is not dead. It is everywhere.}
